\section{Complex-gradient descent with a strict descent margin}
\label{sec:complex-descent}

Theorem 3 in the published version of \citet{zhang} motivates complex-gradient optimization,
but a stationary-point conclusion requires a strict descent margin. The
following convention-explicit statement also applies to a smooth CHNN
regression objective.

\begin{proposition}[Complex descent]
Let $\mathcal L:\mathbb C^P\simeq\mathbb R^{2P}\to\mathbb R$ be bounded
below and have an $L$-Lipschitz real gradient on every segment traversed by
the iterates. For $w=a+ib$, define
$G(w)=\nabla_a\mathcal L+i\nabla_b\mathcal L
=2\partial\mathcal L/\partial\bar w$.
Choose a constant $\rho\in\mathbb C$ with
$\kappa=\operatorname{Re}\rho-(L/2)|\rho|^2>0$, and update
$w_{t+1}=w_t-\rho G(w_t)$.
Then
\[
 \mathcal L(w_{t+1})\leq\mathcal L(w_t)-\kappa\|G(w_t)\|_2^2,
 \qquad \lim_{t\to\infty}\|G(w_t)\|_2=0,
\]
and
\[
 \min_{0\leq t<T}\|G(w_t)\|_2^2
 \leq\frac{\mathcal L(w_0)-\inf\mathcal L}{\kappa T}.
\]
\end{proposition}

\begin{proof}
Under realification, multiplication by $\rho=u+iv$ is $uI+vJ_c$, with
$J_c^\top=-J_c$ and $J_c^\top J_c=I$.
Hence the update has inner product $-u\|G\|^2$ with the real gradient
and squared norm $|\rho|^2\|G\|^2$.
The real descent lemma gives the first inequality. Summing it and using
the lower bound for $\mathcal L$ gives
$\kappa\sum_{t\geq0}\|G(w_t)\|^2\leq\mathcal L(w_0)-\inf\mathcal L$.
Square summability proves convergence of the gradient norm and the minimum
bound. The result asserts convergence to stationarity in gradient norm,
not convergence to a global minimum or to a particular parameter vector.
\end{proof}

For a real step this condition is $0<\rho<2/L$ under the stated convention.
The case $\rho=0$, allowed by a non-strict printed step-size interval,
cannot imply gradient convergence. Nonsmooth complex activations also need
separate treatment; a smooth polynomial CHNN satisfies the differentiability
requirements, with a Lipschitz gradient bound on any bounded iterate region.
This proposition concerns smooth gradient descent; the experimental
readouts are fitted by linear least squares.
